%% Question bank chapter file - generated by bank_tools.py build (do not edit while the bank is still being built:
%% it is regenerated from the per-test data in Question Bank\data). Every question is one bankq block:
%%   \begin{bankq}[id=..., source={...}, paper=.., qno=.., type=..., subtopic={...}, answer={...}, answer_src=key|solved,
%%                 working={...}, difficulty=1-10, diff_parts={K S I C F}, diff_note={...},
%%                 group=<passage id or empty>, also_in={duplicate copies in other tests}]
%% difficulty = min(10, K+S+I+C+F-1): K concepts 1-3, S steps 1-4, I insight 0-2, C computation 0-2, F format 0-1
%%   ...question LaTeX, options, figures (\bfig / \bfigi -> img/)...
%%   \end{bankq}
%% Passages shared by several questions are bankpassage blocks; their questions carry group=<passage id>.
\documentclass[10pt]{article}
\usepackage[a4paper,margin=1.8cm]{geometry}
\usepackage[T1]{fontenc}
\usepackage{amsmath,amssymb,mathtools}
\usepackage{graphicx}
\usepackage[export]{adjustbox}
\usepackage{xcolor,needspace,keyval}
\usepackage{enumitem}
\usepackage{booktabs,array,tabularx}
\usepackage[hidelinks]{hyperref}
\setlength{\parindent}{0pt}
\setlength{\parskip}{3pt}
\makeatletter
\define@key{bq}{id}{\def\bq@id{#1}}\define@key{bq}{source}{\def\bq@src{#1}}\define@key{bq}{paper}{\def\bq@paper{#1}}
\define@key{bq}{qno}{\def\bq@qno{#1}}\define@key{bq}{type}{\def\bq@type{#1}}\define@key{bq}{subtopic}{\def\bq@sub{#1}}
\define@key{bq}{answer}{\def\bq@ans{#1}}\define@key{bq}{answer_src}{\def\bq@asrc{#1}}\define@key{bq}{working}{\def\bq@work{#1}}
\define@key{bq}{group}{\def\bq@group{#1}}\define@key{bq}{also_in}{\def\bq@also{#1}}
\define@key{bq}{difficulty}{\def\bq@diff{#1}}\define@key{bq}{diff_parts}{\def\bq@dparts{#1}}\define@key{bq}{diff_note}{\def\bq@dnote{#1}}
\newenvironment{bankq}[1][]{\def\bq@work{}\def\bq@also{}\def\bq@group{}\def\bq@diff{?}\def\bq@dparts{}\def\bq@dnote{}\setkeys{bq}{#1}%
  \par\needspace{4\baselineskip}\medskip\noindent{\small\textbf{\bq@id}\hfill\color{gray}[\bq@src, Paper \bq@paper, Q\bq@qno]
  \bq@type\ $\cdot$ \bq@sub\ $\cdot$ \textbf{D\bq@diff}}\par\nopagebreak\relax}
  {\par{\small\color{teal}Answer: \textbf{\bq@ans}\ifx\bq@work\@empty\else\ (solved: \bq@work)\fi
  \ifx\bq@also\@empty\else\par\color{gray}Also in: \bq@also\fi}\par}
\define@key{bp}{id}{\def\bp@id{#1}}\define@key{bp}{source}{\def\bp@src{#1}}\define@key{bp}{questions}{\def\bp@qs{#1}}
\newenvironment{bankpassage}[1][]{\setkeys{bp}{#1}\par\needspace{4\baselineskip}\bigskip\noindent
  {\small\textbf{Passage \bp@id}\hfill\color{gray}[\bp@src; questions \bp@qs]}\par\nopagebreak\relax\itshape}{\par}
\makeatother
\newcommand{\bfig}[2]{\par\smallskip{\centering\includegraphics[width=#1pt,max width=0.62\linewidth,max height=0.32\textheight]{img/#2}\par}\smallskip}
\newcommand{\bfigi}[2]{\includegraphics[width=#1pt,max width=\linewidth,max height=0.32\textheight,valign=c]{img/#2}}
\newcommand{\bmissing}[1]{\textit{[figure missing in the source conversion: \detokenize{#1}]}}
\newenvironment{qoptions}{\begin{itemize}[leftmargin=1.8em,itemsep=1pt,topsep=2pt,label={}]}{\end{itemize}}
\newcommand{\opt}[2]{\item[(#1)] #2}
\providecommand{\fwstart}{}\providecommand{\fwend}{}\providecommand{\qfullsplit}{}
\newcolumntype{L}[1]{>{\raggedright\arraybackslash}p{#1}}
\newcolumntype{C}[1]{>{\centering\arraybackslash}p{#1}}
\newcolumntype{R}[1]{>{\raggedleft\arraybackslash}p{#1}}
\providecommand{\sectionhdr}[1]{}\providecommand{\subsectionhdr}[1]{}


\begin{document}
{\Large\bfseries SOT --- Question Bank}\par\smallskip
{\small Mathematics chapter M13. Questions from all converted test papers, grouped by source series. Answers: printed key where it matches the paper, otherwise solved (answer\_src=solved, with working). Exact repeats across papers appear once, with also\_in listing the other copies.}\par\medskip\hrule

\section*{Score-I papers}
\begin{bankq}[id=SI12-P2-MATH-S3-Q08, source={Score-I - Test 12 (09 Mar 2022)}, paper=2, qno=8, type=matching, subtopic={Sine and cosine rules}, answer={P$\to$3, Q$\to$2, R$\to$1}, answer_src=key, difficulty=6, diff_parts={K3 S2 I1 C0 F1}, diff_note={concepts: triangles with consecutive integer sides, complex-number arguments, vector products | steps: P: arg condition gives angle R = 2 angle Q so sides 4,5,6; Q: right angle with sides 3,4,5; R: line conditions; match | I: relate arguments to angles}]
In a $\Delta PQR$, the side lengths $QR$, $RP$ and $PQ$ are consecutive positive integers in increasing order.
\begin{center}\footnotesize
\begin{tabular}{|p{0.6cm}|p{5.3cm}|p{0.6cm}|p{1.0cm}|}
\hline
 & \centering List-I & & \centering List-II\tabularnewline
\hline
(P) & If $z_1$, $z_2$ and $z_3$ are vertices of $P$, $Q$ and $R$ respectively in argand plane, such that $\left|\arg\left(\dfrac{z_1-z_3}{z_2-z_3}\right)\right|=2\left|\arg\left(\dfrac{z_3-z_1}{z_2-z_1}\right)\right|$, then biggest side of triangle is & (1) & 21\\
\hline
(Q) & Let $\vec p$, $\vec q$ and $\vec r$ be the position vectors of vertices $P$, $Q$ and $R$ respectively. If $(\vec r-\vec p)\cdot(\vec q-\vec r)=0$, then the value of $|\vec p\times\vec q+\vec q\times\vec r+\vec r\times\vec p|$ is equal to & (2) & 12\\
\hline
(R) & Let the equations $a_1x+b_1y+c_1=0$ and $a_2x+b_2y+c_2=0$ represent the lines $PQ$ and $PR$ respectively and $\left|\dfrac{a_1b_2-a_2b_1}{a_1a_2+b_1b_2}\right|=\dfrac43$, then the value of the semiperimeter of $\Delta PQR$ is & (3) & 6\\
\hline
\end{tabular}
\end{center}
\end{bankq}

\section*{Score-II papers}
\begin{bankq}[id=SII02-P1-MATH-S2-Q10, source={Score-II - Test 02 (28 Mar 2021)}, paper=1, qno=10, type=multi, subtopic={Sine and cosine rules}, answer={B,C,D}, answer_src=key, difficulty=8, diff_parts={K3 S1 I2 C2 F1}, diff_note={concepts: standard problem in this chapter, sine and cosine rules, apply the relevant results | steps: apply the relevant results; compute}]
Suppose x, y, z are the length of side of a triangle opposite to it angles X, Y and Z respectively if
\[\tan X+\tan Y=\dfrac{-2}{\sqrt3},\quad \tan Y+\tan Z=\dfrac{2}{\sqrt3}\ \text{and}\ x=2\sqrt3\]
then which of the following statement is/are true?

\begin{qoptions}
\opt{A}{X = Y + Z}
\opt{B}{X = 2(Y + Z)}
\opt{C}{$\tan\dfrac{X}{2}=\dfrac{2x}{y+z}$}
\opt{D}{In radius of $\triangle XYZ$ is $\dfrac{\sqrt3}{\left(2+\sqrt3\right)}$}
\end{qoptions}
\end{bankq}
\begin{bankq}[id=SII04-P2-MATH-S4-Q01, source={Score-II - Test 04 (07 May 2021)}, paper=2, qno=1, type=numerical, subtopic={Sine and cosine rules}, answer={7}, answer_src=key, difficulty=8, diff_parts={K3 S1 I2 C2 F1}, diff_note={concepts: standard problem in this chapter, sine and cosine rules, apply the relevant results | steps: apply the relevant results; compute}]
A car is moving at uniform speed towards a tower. It takes 15 minutes for the angle of depression from the top of tower to the car to change from $30^\circ$ to $60^\circ$. The car will reach the base at tower after P minutes. Value of [P] is (Where [.] denotes greatest integer function)
\end{bankq}
\begin{bankq}[id=SII29-P1-MATH-S4-Q01, source={Score-II - Test 29 (26 Mar 2024)}, paper=1, qno=1, type=numerical, subtopic={Sine and cosine rules}, answer={5}, answer_src=key, difficulty=8, diff_parts={K3 S1 I2 C2 F1}, diff_note={concepts: standard problem in this chapter, sine and cosine rules, apply the relevant results | steps: apply the relevant results; compute}]
If the two lines AB : $\left(\displaystyle\int_0^{2t}\left(\dfrac{\sin x}{x}+1\right)dx\right)x+y=3t$ and AC : $2tx+y=0$ intersect at a point A, then x-coordinate of point A as $t\to0$, is equal to $\dfrac pq$ (p and q are in their lowest form). Then (p + q) is
\end{bankq}

\section*{Score Advanced papers}
\begin{bankq}[id=SA01-P1-MATH-S1-Q05, source={Score Advanced - Test 01 (22 Sep 2020)}, paper=1, qno=5, type=single, subtopic={Sine and cosine rules}, answer={D}, answer_src=key, difficulty=6, diff_parts={K2 S2 I2 C1 F0}, diff_note={concepts: ambiguous case of the triangle, sine and cosine rules | steps: quadratic in b; roots relation; expression}]
In $\Delta$ABC with usual notation, if a,c, A are given and b$_1$ = 2b$_2$ where b$_1$, b$_2$ are two values of the third side then the value of $\dfrac{c\sqrt{1+8\sin^2A}}{a}$ is less than
\begin{qoptions}
\opt{A}{1}
\opt{B}{2}
\opt{C}{3}
\opt{D}{4}
\end{qoptions}
\end{bankq}
\begin{bankq}[id=SA01-P2-MATH-S1-Q01, source={Score Advanced - Test 01 (22 Sep 2020)}, paper=2, qno=1, type=multi, subtopic={Circles of a triangle}, answer={A,B,C}, answer_src=key, difficulty=8, diff_parts={K2 S2 I2 C2 F1}, diff_note={concepts: circles on altitudes of triangles, circles of a triangle | steps: circumcentre geometry; radii; relations}]
Let $O$ be the circumcentre of acute angled $\Delta ABC$ and let $r_1,r_2$ and $r_3$ be the radii of the circles drawn on the altitudes $OD$, $OE$ and $OF$ of triangle $OBC$, $OCA$ and $OAB$ as diameter, then minimum value of $\dfrac{a^2}{r_1^2}+\dfrac{b^2}{r_2^2}+\dfrac{c^2}{r_3^2}$ is divisible by :
\begin{qoptions}
\opt{A}{36}
\opt{B}{72}
\opt{C}{144}
\opt{D}{216}
\end{qoptions}
\end{bankq}
\begin{bankq}[id=SA02-P1-MATH-S3-Q01, source={Score Advanced - Test 02 (13 May 2021)}, paper=1, qno=1, type=numerical, subtopic={Sine and cosine rules}, answer={1.2}, answer_src=solved, working={tan((A-B)/2)=1/sqrt63 gives cos C=1/8 and c=6}, difficulty=7, diff_parts={K2 S2 I2 C1 F1}, diff_note={concepts: triangle from sides and an angle difference, sine and cosine rules | steps: law of tangents; half angle; third side}]
In a triangle ABC (with usual notations), a = 5, b = 4 and A $-$ B = $\cos^{-1}\left(\dfrac{31}{32}\right)$ then the value of $\dfrac{c}{5}$ is
\end{bankq}
\begin{bankq}[id=SA02-P2-MATH-S3-Q08, source={Score Advanced - Test 02 (13 May 2021)}, paper=2, qno=8, type=single, subtopic={Circles of a triangle}, answer={1}, answer_src=solved, working={QR=k; area and sum of squares conditions give k=1}, difficulty=7, diff_parts={K2 S2 I2 C2 F0}, diff_note={concepts: triangle in a circle, circles of a triangle | steps: chord and angle; area formula; sum of squares; radius}]
Let a circle with centre O and radius k. 3 points P, Q, R are chosen on a circle such that if length of chord PQ, QR, RP are x, y \& z respectively, x$^2$ + y$^2$ + z$^2$ = 5, $\angle$QOR = 60$^\circ$ and Area of $\Delta$PQR is $\dfrac{\sqrt{3}}{4}$ unit, then the value of radius of circle k will be equal to
\end{bankq}
\begin{bankq}[id=SA05-P1-MATH-S2-Q05, source={Score Advanced - Test 05 (15 Jul 2021)}, paper=1, qno=5, type=numerical, subtopic={Sine and cosine rules}, answer={34.64}, answer_src=key, difficulty=8, diff_parts={K3 S2 I2 C1 F1}, diff_note={concepts: towers with depression and elevation angles, sine and cosine rules, tan(theta)=60/d | steps: tan(theta)=60/d; tan(phi)=20/d; theta=2 phi}]
A tower T$_1$ of height 60m is located exactly opposite to tower T$_2$ of height 80m on a straight road. From top of T$_1$, if angle of depression of foot of T$_2$ is twice the angle of elevation of top of T$_2$, then distance (in meter) between tower T$_1$ and tower T$_2$ is-
\end{bankq}
\begin{bankq}[id=SA06-P1-MATH-S1-Q05, source={Score Advanced - Test 06 (04 Sep 2021)}, paper=1, qno=5, type=multi, subtopic={Sine and cosine rules}, answer={A,B,C}, answer_src=key, difficulty=8, diff_parts={K3 S1 I2 C2 F1}, diff_note={concepts: triangle from inverse trigonometric angles, sine and cosine rules, tan B=2, tan C=3, tan A=1 | steps: tan B=2, tan C=3, tan A=1; sides from the sine rule}]
In $\Delta$ABC, if $\angle$B = $\sec^{-1}\left(\dfrac{5}{4}\right)$ + $\text{cosec}^{-1}\sqrt{5}$, $\angle$C = $\text{cosec}^{-1}\left(\dfrac{25}{7}\right)$ + $\cot^{-1}\left(\dfrac{9}{13}\right)$ and c = 3. Which of the following statement(s) is/are correct ? (where all symbols used have their usual meanings in a triangle)
\begin{qoptions}
\opt{A}{tan A, tan B, tan C are in arithmetic progression.}
\opt{B}{Inradius the triangle is irrational}
\opt{C}{Area of $\Delta$ABC is rational.}
\opt{D}{Semiperimeter is rational}
\end{qoptions}
\end{bankq}
\begin{bankq}[id=SA06-P2-MATH-S1-Q06, source={Score Advanced - Test 06 (04 Sep 2021)}, paper=2, qno=6, type=multi, subtopic={Sine and cosine rules}, answer={B,C}, answer_src=key, difficulty=8, diff_parts={K3 S1 I2 C2 F1}, diff_note={concepts: triangle relation from cos A + cos B, sine and cosine rules, use half-angle formulas | steps: use half-angle formulas; cot(A/2)cot(B/2)=3}]
In $\Delta$ABC (with usual notation), if cos A + cos B = 4sin$^2\dfrac{C}{2}$, then which of the following hold(s) good ?
\begin{qoptions}
\opt{A}{cot$\dfrac{A}{2}$cot$\dfrac{B}{2}$ = 2}
\opt{B}{cot$\dfrac{A}{2}$cot$\dfrac{B}{2}$ = 3}
\opt{C}{a, c, b are in A.P.}
\opt{D}{a, b, c are in G.P.}
\end{qoptions}
\end{bankq}
\begin{bankq}[id=SA11-P1-MATH-S1-Q06, source={Score Advanced - Test 11 (09 May 2023)}, paper=1, qno=6, type=multi, subtopic={Circles of a triangle}, answer={A,C}, answer_src=key, difficulty=8, diff_parts={K3 S1 I2 C2 F1}, diff_note={concepts: sequence of excentral triangles, circles of a triangle, angles of the excentral triangle are (pi-A)/2 etc. | steps: angles of the excentral triangle are (pi-A)/2 etc.}, also_in={Score Advanced - Test 29 (01 May 2024) P1 Q6}]
For $\Delta$ABC, let A$_1$B$_1$C$_1$ be the triangle formed by the excentres of $\Delta$ABC, let $\Delta$A$_2$B$_2$C$_2$ be formed by the excentres of $\Delta$A$_1$B$_1$C$_1$ and so on. Finally let $\Delta$A$_n$B$_n$C$_n$ be the n$^\text{th}$ derived triangle that it is formed by the excentres of $\Delta$A$_{n-1}$ B$_{n-1}$C$_{n-1}$ then
\begin{qoptions}
\opt{A}{$\angle$A$_1$ = 90\textdegree{} $-$ $\dfrac{A}{2}$}
\opt{B}{The circumcentre of $\Delta$ABC always coincides with the orthocentre of $\Delta$A$_1$B$_1$C$_1$}
\opt{C}{as n$\longrightarrow\infty$, $\Delta$A$_n$B$_n$C$_n$ tends to become equilateral}
\opt{D}{B$_1$C$_1$ = a cot $\dfrac{A}{2}$}
\end{qoptions}
\end{bankq}
\begin{bankq}[id=SA11-P1-MATH-S2-Q07, source={Score Advanced - Test 11 (09 May 2023)}, paper=1, qno=7, type=single, subtopic={Circles of a triangle}, answer={A}, answer_src=key, difficulty=7, diff_parts={K3 S1 I2 C2 F0}, diff_note={concepts: distances between triangle centres, circles of a triangle, Euler line relations OG:GH=1:2 | steps: Euler line relations OG:GH=1:2}]
Let O, O$^1$ and G be the circumcentre, orthocentre and centroid of $\Delta$ABC respectively. Now match the values in list-I and list-II.\par
{\footnotesize\setlength{\tabcolsep}{1.5pt}\noindent\begin{tabular}{|c|L{4.6cm}|c|c|}
\hline
 & \multicolumn{1}{c|}{list-I} & & \multicolumn{1}{c|}{list-II}\\ \hline
(A) & If S is any point in plane of A,B,C then $\dfrac{\left|\overrightarrow{\text{SA}} + \overrightarrow{\text{SB}} + \overrightarrow{\text{SC}}\right|}{\left|\overrightarrow{\text{SG}}\right|}$ = & (P) & 4\\ \hline
(B) & $\dfrac{\left|\overrightarrow{\text{OA}} + \overrightarrow{\text{OB}} + \overrightarrow{\text{OC}}\right|}{\left|\overrightarrow{\text{OO}'}\right|}$ = & (Q) & 3\\ \hline
(C) & $\dfrac{\left|\overrightarrow{\text{O}'\text{A}} + \overrightarrow{\text{O}'\text{B}} + \overrightarrow{\text{O}'\text{C}}\right|}{\left|\overrightarrow{\text{OO}'}\right|}$ = & (R) & 2\\ \hline
(D) & $\dfrac{\left|\overrightarrow{\text{AO}'} + \overrightarrow{\text{O}'\text{B}} + \overrightarrow{\text{O}'\text{C}}\right|}{\left|\overrightarrow{\text{AP}}\right|}$ = & (S) & 1\\ \hline
 & & (T) & 1/2\\ \hline
\end{tabular}}\par
(Where $\left|\overrightarrow{\text{AP}}\right|$ is diameter of circumcircle of $\Delta$ABC)
\begin{qoptions}
\opt{A}{A $\longrightarrow$ Q;B $\longrightarrow$ S;C $\longrightarrow$ R;D $\longrightarrow$ S}
\opt{B}{A $\longrightarrow$ Q;B $\longrightarrow$ R;C $\longrightarrow$ S;D $\longrightarrow$ S}
\opt{C}{A $\longrightarrow$ S;B $\longrightarrow$ S;C $\longrightarrow$ Q;D $\longrightarrow$ R}
\opt{D}{A $\longrightarrow$ S;B $\longrightarrow$ Q;C $\longrightarrow$ R;D $\longrightarrow$ R}
\end{qoptions}
\end{bankq}
\begin{bankq}[id=SA14-P1-MATH-S3-Q02, source={Score Advanced - Test 14 (16 May 2023)}, paper=1, qno=2, type=numerical, subtopic={Circles of a triangle}, answer={17.125}, answer_src=key, difficulty=8, diff_parts={K3 S1 I2 C2 F1}, diff_note={concepts: triangle with a given incircle and tangency lengths, circles of a triangle, area from r and s | steps: area from r and s; circumradius}]
Let incircle of radius 4 units of a triangle ABC touches the side BC at D. If BD = 6, DC = 8 and $\Delta$ be the area of triangle, and R be radius of circumcircle of $\Delta$ABC, then $\sqrt{\Delta-3}$ + R is equal to \_\_\_\_\_\_\_\_\_\_.\par
\end{bankq}
\begin{bankq}[id=SA14-P1-MATH-S3-Q03, source={Score Advanced - Test 14 (16 May 2023)}, paper=1, qno=3, type=numerical, subtopic={Sine and cosine rules}, answer={3}, answer_src=solved, working={At theta=60 degrees both sides equal 1.577, so k=3}, difficulty=8, diff_parts={K3 S1 I2 C2 F1}, diff_note={concepts: tangential quadrilateral inside a semicircle, sine and cosine rules, OB+CD=OD+BC with chord lengths 2cos(theta), 2sin(theta)-1/sin(theta), cot(theta) | steps: OB+CD=OD+BC with chord lengths 2cos(theta), 2sin(theta)-1/sin(theta), cot(theta)}, also_in={Score Advanced - Test 20 (06 Apr 2024) P1 Q3}]
On a semicircle with diameter AB and centre O point E and C are marked in such a way that OE is perpendicular to AB, and the chord AC intersects the segment OE at a point D which is interior to the semicircle. If $\angle$CBA = $\dfrac{\pi}{\text{k}}$ such that a circle can be inscribed into the quadrilateral OBCD then value of k is equal to\par
\end{bankq}
\begin{bankq}[id=SA20-P1-MATH-S3-Q02, source={Score Advanced - Test 20 (06 Apr 2024)}, paper=1, qno=2, type=numerical, subtopic={Circles of a triangle}, answer={17.125}, answer_src=solved, working={With r=4, s-b=6 and s-c=8 and s-a=x, area = r s = sqrt(s x 48) gives x=7, s=21, delta=84. Sides a=14, b=15, c=13 and R=abc/(4 delta)=8.125. sqrt(84-3)+8.125 = 17.125.}, difficulty=9, diff_parts={K3 S2 I2 C2 F1}, diff_note={concepts: area and circumradius from incircle tangent lengths, circles of a triangle, s-a from the inradius | steps: s-a from the inradius; area; sides; R=abc/4 delta}]
Let incircle of radius 4 units of a triangle ABC touches the side BC at D. If BD $=6$, DC $=8$ and $\Delta$ be the area of triangle, and $R$ be radius of circumcircle of $\Delta$ABC, then $\sqrt{\Delta-3}+R$ is equal to \underline{\hspace{1.5cm}}.
\end{bankq}
\begin{bankq}[id=SA24-P2-MATH-S1-Q02, source={Score Advanced - Test 24 (17 Apr 2024)}, paper=2, qno=2, type=single, subtopic={Area}, answer={B}, answer_src=key, difficulty=7, diff_parts={K3 S1 I2 C2 F0}, diff_note={concepts: triangle quantity from segments of the base and tan A, area, tan A from tan B and tan C | steps: tan A from tan B and tan C}]
In a acute angle triangle ABC, the foot of the perpendicular from A divide the opposite side into two parts of length 3 and 17 units and $\tan A=\dfrac{22}{7}$, the $\left(\dfrac{\text{Area of }\Delta ABC}{22}\right)$ is
\begin{qoptions}
\opt{A}{$\dfrac52$}
\opt{B}{5}
\opt{C}{10}
\opt{D}{20}
\end{qoptions}
\end{bankq}
\begin{bankq}[id=SA24-P2-MATH-S1-Q04, source={Score Advanced - Test 24 (17 Apr 2024)}, paper=2, qno=4, type=single, subtopic={Sine and cosine rules}, answer={C}, answer_src=key, difficulty=7, diff_parts={K3 S1 I2 C2 F0}, diff_note={concepts: ratio of heights seen at equal angles from a point, sine and cosine rules, angles subtended by the pillar and flagstaff | steps: angles subtended by the pillar and flagstaff}]
A flagstaff stands vertically on a pillar, the height of the flagstaff being double the height of the pillar. A man on the ground at a distance finds that both the pillar and the flagstaff subtend equal angles at his eyes. The ratio of the height of the pillar and the distance of the man from the pillar, is
\begin{qoptions}
\opt{A}{$\sqrt3:1$}
\opt{B}{1 : 3}
\opt{C}{$1:\sqrt3$}
\opt{D}{$\sqrt3:2$}
\end{qoptions}
\end{bankq}
\begin{bankq}[id=SA24-P2-MATH-S2-Q10, source={Score Advanced - Test 24 (17 Apr 2024)}, paper=2, qno=10, type=multi, subtopic={Sine and cosine rules}, answer={B,D}, answer_src=key, difficulty=8, diff_parts={K3 S1 I2 C2 F1}, diff_note={concepts: angles impossible in a triangle with sides in GP, sine and cosine rules, cosine rule with b\^{}2=ac | steps: cosine rule with b\^{}2=ac}]
If the sides a, b \& c of a triangle ABC are in G.P. then $\angle B$ can never be
\begin{qoptions}
\opt{A}{$\dfrac\pi6$}
\opt{B}{$\dfrac{5\pi}{6}$}
\opt{C}{$\dfrac\pi4$}
\opt{D}{$\dfrac{2\pi}{3}$}
\end{qoptions}
\end{bankq}
\begin{bankq}[id=SA29-P1-MATH-S2-Q07, source={Score Advanced - Test 29 (01 May 2024)}, paper=1, qno=7, type=single, subtopic={Circles of a triangle}, answer={A}, answer_src=key, difficulty=7, diff_parts={K3 S1 I2 C2 F0}, diff_note={concepts: duplicate or standard problem in this chapter, circles of a triangle, apply the relevant results | steps: apply the relevant results; compute}]
Let O, O$^1$ and G be the circumcentre, orthocentre and centroid of $\Delta$ABC respectively. Now match the values in list-I and list-II.\par
{\footnotesize\setlength{\tabcolsep}{1.5pt}\noindent\begin{tabular}{|c|L{5.2cm}|c|L{0.8cm}|}
\hline
 & \multicolumn{1}{>{\centering\arraybackslash}p{5.2cm}|}{List-I} & & \multicolumn{1}{>{\centering\arraybackslash}p{0.8cm}|}{List-II} \\ \hline
(A) & If S is any point in plane of A,B,C then $\dfrac{\left|\overrightarrow{SA}+\overrightarrow{SB}+\overrightarrow{SC}\right|}{\left|\overrightarrow{SG}\right|}$ = & (P) & 4 \\ \hline
(B) & $\dfrac{\left|\overrightarrow{OA}+\overrightarrow{OB}+\overrightarrow{OC}\right|}{\left|\overrightarrow{OO'}\right|}$ = & (Q) & 3 \\ \hline
(C) & $\dfrac{\left|\overrightarrow{O'A}+\overrightarrow{O'B}+\overrightarrow{O'C}\right|}{\left|\overrightarrow{OO'}\right|}$ = & (R) & 2 \\ \hline
(D) & $\dfrac{\left|\overrightarrow{AO'}+\overrightarrow{O'B}+\overrightarrow{O'C}\right|}{\left|\overrightarrow{AP}\right|}$ = & (S) & 1 \\ \hline
 & & (T) & 1/2 \\ \hline
\end{tabular}}\par
(Where $\left|\overrightarrow{AP}\right|$ is diameter of circumcircle of $\Delta$ABC)
\begin{qoptions}
\opt{A}{A $\rightarrow$ Q;B $\rightarrow$ S;C $\rightarrow$ R;D $\rightarrow$ S}
\opt{B}{A $\rightarrow$ Q;B $\rightarrow$ R;C $\rightarrow$ S;D $\rightarrow$ S}
\opt{C}{A $\rightarrow$ S;B $\rightarrow$ S;C $\rightarrow$ Q;D $\rightarrow$ R}
\opt{D}{A $\rightarrow$ S;B $\rightarrow$ Q;C $\rightarrow$ R;D $\rightarrow$ R}
\end{qoptions}
\end{bankq}
\begin{bankq}[id=SA33-P1-MATH-S4-Q06, source={Score Advanced - Test 33 (12 May 2024)}, paper=1, qno=6, type=numerical, subtopic={Area}, answer={23}, answer_src=key, difficulty=8, diff_parts={K3 S1 I2 C2 F1}, diff_note={concepts: standard problem in this chapter, area, apply the relevant results | steps: apply the relevant results; compute}, also_in={Score Advanced - Test 37 (19 May 2024) P1 Q6}]
Let the area of given triangle ABC be $\Delta$. Points A$_1$, B$_1$ and C$_1$ are the mid points of the BC, CA and AB respectively. Point A$_2$ is the mid point of CA$_1$. Lines C$_1$A$_1$ and AA$_2$ meet the median BB$_1$ at points E and D respectively. If $\Delta_1$ be the area of the quadrilateral A$_1$A$_2$DE. If the value of $\dfrac{\Delta_1}\Delta=\dfrac mn$ (where m and n are coprime), then n $-$ 3m = \_\_\_\_\_.
\end{bankq}
\begin{bankq}[id=SA40-P2-MATH-S1-Q01, source={Score Advanced - Test 40 (27 Apr 2025)}, paper=2, qno=1, type=multi, subtopic={Sine and cosine rules}, answer={A,B,D}, answer_src=key, difficulty=8, diff_parts={K3 S1 I2 C2 F1}, diff_note={concepts: standard problem in this chapter, sine and cosine rules, apply the relevant results | steps: apply the relevant results; compute}]
Which of the following is/are true.
\begin{qoptions}
\opt{A}{In any $\Delta$ABC, min value of\par{\small$\displaystyle\sum\dfrac{\sqrt{\sin A}}{\sqrt{\sin B}+\sqrt{\sin C}-\sqrt{\sin A}}$}\par is 3.}
\opt{B}{If $\sec A\tan B+\tan A\sec B=30$ then $\sec A\sec B+\tan A\tan B=\sqrt{901}$.}
\opt{C}{If $\sec A\tan B+\tan A\sec B=30$ then $\sec A\sec B+\tan A\tan B=\sqrt{899}$.}
\opt{D}{Three lines $y-z-1=0=x$, $z+x+1=0=y$, $x-z-1=0=y$ intersect the xy plane at U, M, R. Then the orthocenter of triangle UMR is $(0,1,0)$.}
\end{qoptions}
\end{bankq}
\begin{bankq}[id=SA40-P2-MATH-S4-Q05, source={Score Advanced - Test 40 (27 Apr 2025)}, paper=2, qno=5, type=numerical, subtopic={Circles of a triangle}, answer={6}, answer_src=key, difficulty=8, diff_parts={K3 S1 I2 C2 F1}, diff_note={concepts: standard problem in this chapter, circles of a triangle, apply the relevant results | steps: apply the relevant results; compute}]
On the side AC of an acute triangle ABC a point D is taken such that AD = 1, DC = 2 and BD is an altitude of the triangle. A circle of radius 2 passes through the points A and D and touches at the point D the circle circumscribed about the triangle BDC. The area of $\Delta$ABC is $p\sqrt q$ ($p$ is prime and $q$ is natural number). Then the value of $\dfrac{q-p}2$ is :
\end{bankq}

\section*{SRGBT papers}
\begin{bankq}[id=SRGBT03-P1-MATH-S1-Q41, source={SRGBT - Test 03 (08 Aug 2021)}, paper=1, qno=41, type=multi, subtopic={Circles of a triangle}, answer={A,C}, answer_src=key, difficulty=8, diff_parts={K3 S1 I2 C2 F1}, diff_note={concepts: standard problem in this chapter, circles of a triangle, apply the relevant results | steps: apply the relevant results; compute}]
Three equal circles of radius 1 touch one another. Radius of the circle touching all the three circles can be :
\begin{qoptions}
\opt{A}{$\dfrac{2-\sqrt3}{\sqrt3}$}
\opt{B}{$\dfrac{\sqrt3-\sqrt2}{\sqrt2}$}
\opt{C}{$\dfrac{2+\sqrt3}{\sqrt3}$}
\opt{D}{$\dfrac{\sqrt3+\sqrt2}{\sqrt2}$}
\end{qoptions}
\end{bankq}
\begin{bankq}[id=SRGBT03-P1-MATH-S1-Q43, source={SRGBT - Test 03 (08 Aug 2021)}, paper=1, qno=43, type=multi, subtopic={Sine and cosine rules}, answer={A,C,D}, answer_src=key, difficulty=8, diff_parts={K3 S1 I2 C2 F1}, diff_note={concepts: standard problem in this chapter, sine and cosine rules, apply the relevant results | steps: apply the relevant results; compute}]
The sides of a $\Delta$ABC satisfy the equation, 2a$^2$ + 4b$^2$ + c$^2$ = 4ab + 2ac. Then :
\begin{qoptions}
\opt{A}{the triangle is isosceles}
\opt{B}{the triangle is obtuse}
\opt{C}{B = cos$^{-1}\dfrac78$}
\opt{D}{A = cos$^{-1}\dfrac14$}
\end{qoptions}
\end{bankq}
\begin{bankq}[id=SRGBT03-P1-MATH-S1-Q44, source={SRGBT - Test 03 (08 Aug 2021)}, paper=1, qno=44, type=multi, subtopic={Sine and cosine rules}, answer={A,B,D}, answer_src=key, difficulty=8, diff_parts={K3 S1 I2 C2 F1}, diff_note={concepts: standard problem in this chapter, sine and cosine rules, apply the relevant results | steps: apply the relevant results; compute}]
If in a $\Delta$ABC, cosA . cosB + sinA sinB sin2C = 1 then,
\begin{qoptions}
\opt{A}{$\Delta$ABC is isosceles but not right angled}
\opt{B}{$\Delta$ABC is acute angled}
\opt{C}{$\Delta$ABC is right angled}
\opt{D}{least angle of the triangle is $\dfrac\pi4$}
\end{qoptions}
\end{bankq}
\begin{bankq}[id=SRGBT03-P1-MATH-S1-Q47, source={SRGBT - Test 03 (08 Aug 2021)}, paper=1, qno=47, type=multi, subtopic={Sine and cosine rules}, answer={B,C}, answer_src=key, difficulty=8, diff_parts={K3 S1 I2 C2 F1}, diff_note={concepts: standard problem in this chapter, sine and cosine rules, apply the relevant results | steps: apply the relevant results; compute}]
The sines of two angles of a triangle are equal to $\dfrac5{13}$ and $\dfrac{99}{101}$ . The cosine of the third angle can be :
\begin{qoptions}
\opt{A}{$\dfrac{245}{1313}$}
\opt{B}{$\dfrac{255}{1313}$}
\opt{C}{$\dfrac{735}{1313}$}
\opt{D}{$\dfrac{765}{1313}$}
\end{qoptions}
\end{bankq}
\begin{bankq}[id=SRGBT03-P1-MATH-S1-Q48, source={SRGBT - Test 03 (08 Aug 2021)}, paper=1, qno=48, type=multi, subtopic={Sine and cosine rules}, answer={A,B,C}, answer_src=key, difficulty=8, diff_parts={K3 S1 I2 C2 F1}, diff_note={concepts: standard problem in this chapter, sine and cosine rules, apply the relevant results | steps: apply the relevant results; compute}]
In a $\Delta$ABC, if the median, bisector and altitude drawn from the vertex A divide the angle at the vertex into four equal parts then the angle of the $\Delta$ABC can be equal to :
\begin{qoptions}
\opt{A}{$\dfrac\pi8$}
\opt{B}{$\dfrac\pi2$}
\opt{C}{$\dfrac{3\pi}8$}
\opt{D}{$\dfrac\pi5$}
\end{qoptions}
\end{bankq}
\begin{bankq}[id=SRGBT03-P1-MATH-S1-Q51, source={SRGBT - Test 03 (08 Aug 2021)}, paper=1, qno=51, type=multi, subtopic={Circles of a triangle}, answer={9}, answer_src=key, difficulty=8, diff_parts={K3 S1 I2 C2 F1}, diff_note={concepts: standard problem in this chapter, circles of a triangle, apply the relevant results | steps: apply the relevant results; compute}]
In an isosceles $\Delta$ABC if the altitudes intersect on the inscribed circle then the secant of the vertical angle `A' is :
\end{bankq}
\begin{bankq}[id=SRGBT03-P1-MATH-S1-Q54, source={SRGBT - Test 03 (08 Aug 2021)}, paper=1, qno=54, type=multi, subtopic={Circles of a triangle}, answer={2}, answer_src=key, difficulty=8, diff_parts={K3 S1 I2 C2 F1}, diff_note={concepts: standard problem in this chapter, circles of a triangle, apply the relevant results | steps: apply the relevant results; compute}]
If the circum centre of the $\Delta$ABC lies on its incircle then the value of [R/r] is where [ . ] is GI function :
\end{bankq}
\begin{bankq}[id=SRGBT03-P1-MATH-S1-Q55, source={SRGBT - Test 03 (08 Aug 2021)}, paper=1, qno=55, type=multi, subtopic={Sine and cosine rules}, answer={1}, answer_src=key, difficulty=8, diff_parts={K3 S1 I2 C2 F1}, diff_note={concepts: standard problem in this chapter, sine and cosine rules, apply the relevant results | steps: apply the relevant results; compute}]
Let a, b and c be the sides of a $\Delta$ABC. If a$^2$, b$^2$ and c$^2$ are the roots of the equation x$^3-$8x$^2$+6x $-$ 4 = 0, then the value of $\dfrac{\cos\mathrm A}{\mathrm a}+\dfrac{\cos\mathrm B}{\mathrm b}+\dfrac{\cos\mathrm C}{\mathrm c}$ is equal to :
\end{bankq}
\begin{bankq}[id=SRGBT10-P1-MATH-S1-Q41, source={SRGBT - Test 10 (26 Sep 2021)}, paper=1, qno=41, type=multi, subtopic={Sine and cosine rules}, answer={A}, answer_src=key, difficulty=5, diff_parts={K2 S1 I1 C1 F1}, diff_note={concepts: standard problem in this chapter, sine and cosine rules | steps: apply the relevant results; compute}]
Let a, b, c be the sides of a triangle such that $\dfrac{bc}{b+c}+\dfrac{ca}{c+a}+\dfrac{ab}{a+b}=s$ where s is semi -- perimeter, then the triangle is
\begin{qoptions}
\opt{A}{equilateral}
\opt{B}{isosceles but not right angled}
\opt{C}{right angled but not isosceles}
\opt{D}{obtuse angled}
\end{qoptions}
\end{bankq}
\begin{bankq}[id=SRGBT10-P1-MATH-S1-Q43, source={SRGBT - Test 10 (26 Sep 2021)}, paper=1, qno=43, type=multi, subtopic={Sine and cosine rules}, answer={B,C}, answer_src=key, difficulty=8, diff_parts={K3 S1 I2 C2 F1}, diff_note={concepts: standard problem in this chapter, sine and cosine rules, apply the relevant results | steps: apply the relevant results; compute}]
For a $\triangle$ABC with integer side lengths and $\cos\mathrm A=\dfrac{11}{16}$, $\cos\mathrm B=\dfrac78$, $\cos\mathrm C=-\dfrac14$, which of the following statement(s) is/are true?
\begin{qoptions}
\opt{A}{least possible perimeter of triangle ABC is 27}
\opt{B}{least possible perimeter of triangle ABC is 9}
\opt{C}{least area of the triangle ABC is $\dfrac{3\sqrt{15}}4$}
\opt{D}{None of these}
\end{qoptions}
\end{bankq}
\begin{bankq}[id=SRGBT10-P1-MATH-S1-Q45, source={SRGBT - Test 10 (26 Sep 2021)}, paper=1, qno=45, type=multi, subtopic={Area}, answer={C}, answer_src=key, difficulty=5, diff_parts={K2 S1 I1 C1 F1}, diff_note={concepts: standard problem in this chapter, area | steps: apply the relevant results; compute}]
In $\triangle$ABC, P is any point inside a triangle such that area of $\triangle$BPC, $\triangle$APC, $\triangle$APB are equal. Line AP cut BC at M, area of $\triangle$PMC is 5 sq. unit then area of $\triangle$ABC can be :
\begin{qoptions}
\opt{A}{20 sq. units}
\opt{B}{25 sq. units}
\opt{C}{30 sq. units}
\opt{D}{10 sq. units}
\end{qoptions}
\end{bankq}
\begin{bankq}[id=SRGBT10-P1-MATH-S1-Q46, source={SRGBT - Test 10 (26 Sep 2021)}, paper=1, qno=46, type=multi, subtopic={Sine and cosine rules}, answer={C}, answer_src=key, difficulty=5, diff_parts={K2 S1 I1 C1 F1}, diff_note={concepts: standard problem in this chapter, sine and cosine rules | steps: apply the relevant results; compute}]
In $\triangle$ABC \ $\sin(\mathrm A-\mathrm B)=\dfrac35$, $\sin\mathrm A=2\sin\mathrm B$ \ then \ $\dfrac{\mathrm a+\mathrm b}{\mathrm c}=$
\begin{qoptions}
\opt{A}{$\dfrac4{\sqrt5}$}
\opt{B}{$\sqrt{\dfrac25}$}
\opt{C}{$\dfrac3{\sqrt5}$}
\opt{D}{$\dfrac2{\sqrt5}$}
\end{qoptions}
\end{bankq}
\begin{bankq}[id=SRGBT10-P1-MATH-S1-Q47, source={SRGBT - Test 10 (26 Sep 2021)}, paper=1, qno=47, type=multi, subtopic={Sine and cosine rules}, answer={A,B,C,D}, answer_src=key, difficulty=8, diff_parts={K3 S1 I2 C2 F1}, diff_note={concepts: standard problem in this chapter, sine and cosine rules, apply the relevant results | steps: apply the relevant results; compute}]
In $\triangle$ABC, if a = 2, B = 60$^\circ$, C = 75$^\circ$, then
\begin{qoptions}
\opt{A}{$\mathrm b=\sqrt6$ units}
\opt{B}{area of $\triangle$ABC $=\left(\dfrac{3+\sqrt3}2\right)$ sq. units.}
\opt{C}{Length of angle bisector of $\angle$B is 2 units}
\opt{D}{Circumradius of the triangle ABC is $\sqrt2$ units}
\end{qoptions}
\end{bankq}
\begin{bankq}[id=SRGBT10-P1-MATH-S1-Q49, source={SRGBT - Test 10 (26 Sep 2021)}, paper=1, qno=49, type=multi, subtopic={Sine and cosine rules}, answer={A,C}, answer_src=key, difficulty=8, diff_parts={K3 S1 I2 C2 F1}, diff_note={concepts: standard problem in this chapter, sine and cosine rules, apply the relevant results | steps: apply the relevant results; compute}]
If in a $\triangle$ABC, a, b, c are in A.P. then it is necessary that :
\begin{qoptions}
\opt{A}{$\dfrac23<\dfrac bc<2$}
\opt{B}{$\dfrac13<\dfrac bc<\dfrac23$}
\opt{C}{$\dfrac23<\dfrac ba<2$}
\opt{D}{$\dfrac13<\dfrac ba<\dfrac23$}
\end{qoptions}
\end{bankq}
\begin{bankq}[id=SRGBT10-P1-MATH-S1-Q51, source={SRGBT - Test 10 (26 Sep 2021)}, paper=1, qno=51, type=multi, subtopic={Circles of a triangle}, answer={7}, answer_src=key, difficulty=8, diff_parts={K3 S1 I2 C2 F1}, diff_note={concepts: standard problem in this chapter, circles of a triangle, apply the relevant results | steps: apply the relevant results; compute}]
A triangle ABC has A $-$ B = 120$^\circ$ and the circumradius is 8 times the in--radius. If $\cos\mathrm C=\mathrm a/8$, then a =
\end{bankq}
\begin{bankq}[id=SRGBT10-P1-MATH-S1-Q56, source={SRGBT - Test 10 (26 Sep 2021)}, paper=1, qno=56, type=multi, subtopic={Area}, answer={1}, answer_src=key, difficulty=8, diff_parts={K3 S1 I2 C2 F1}, diff_note={concepts: standard problem in this chapter, area, apply the relevant results | steps: apply the relevant results; compute}]
ABCD is a trapezium. AB and CD are parallel and AD is perpendicular to them. Let AB = 4, CD = 3 and r be the radius of the circle which touches all the sides of trapezium them the value of [r] is (where [.] denotes greatest integer function)
\end{bankq}
\begin{bankq}[id=SRGBT10-P1-MATH-S1-Q58, source={SRGBT - Test 10 (26 Sep 2021)}, paper=1, qno=58, type=multi, subtopic={Sine and cosine rules}, answer={6}, answer_src=key, difficulty=8, diff_parts={K3 S1 I2 C2 F1}, diff_note={concepts: standard problem in this chapter, sine and cosine rules, apply the relevant results | steps: apply the relevant results; compute}]
The sides a, b, c of $\triangle$ABC satisfy the relations $\mathrm a+\mathrm b-\mathrm c=2$ and $2\mathrm{ab}-\mathrm c^2=4$, if $x\in[0,1]$ then max value of $ax^2+bx+c$ is :
\end{bankq}
\begin{bankq}[id=SRGBT10-P1-MATH-S1-Q59, source={SRGBT - Test 10 (26 Sep 2021)}, paper=1, qno=59, type=multi, subtopic={Area}, answer={6}, answer_src=key, difficulty=8, diff_parts={K3 S1 I2 C2 F1}, diff_note={concepts: standard problem in this chapter, area, apply the relevant results | steps: apply the relevant results; compute}]
In parallelogram ABCD, P is any point inside parallelogram such that sum of area of $\triangle$APD and $\triangle$BPC is 3 sq. unit, then area of parallelogram ABCD is
\end{bankq}

\section*{Enthusiast Pre-Score papers}
\begin{bankq}[id=EPS02-P1-MATH-S1-Q01, source={Enthusiast Pre-Score - Test 02 (11 Nov 2020)}, paper=1, qno=1, type=single, subtopic={Sine and cosine rules}, answer={A}, answer_src=key, difficulty=7, diff_parts={K2 S2 I2 C2 F0}, diff_note={concepts: triangle with a point dividing a side, sine and cosine rules | steps: sine rule in two triangles; ratio condition; solve the angle}]
In a triangle ABC, with usual notation $\angle$ABC = 45$^\circ$, Point D is on BC so that 2BD = CD and $\angle$DAB = 75$^\circ$ then $\angle$ACB equals to
\begin{qoptions}
\opt{A}{15$^\circ$}
\opt{B}{$\angle$A $-$ $\angle$B}
\opt{C}{30$^\circ$}
\opt{D}{$\dfrac{\angle A-\angle B}{2}$}
\end{qoptions}
\end{bankq}
\begin{bankq}[id=EPS02-P1-MATH-S1-Q02, source={Enthusiast Pre-Score - Test 02 (11 Nov 2020)}, paper=1, qno=2, type=single, subtopic={Sine and cosine rules}, answer={A}, answer_src=key, difficulty=6, diff_parts={K2 S2 I2 C1 F0}, diff_note={concepts: expression in a triangle with B=2A, sine and cosine rules | steps: relations of sides; substitute; simplify}]
In $\Delta$ABC if B = 2A then $\dfrac{b^2c+a^3}{b^2+c^2}$ is equal to (symbols have their usual meaning)
\begin{qoptions}
\opt{A}{a}
\opt{B}{b}
\opt{C}{c}
\opt{D}{a + b $-$ c}
\end{qoptions}
\end{bankq}
\begin{bankq}[id=EPS02-P1-MATH-S2-Q09, source={Enthusiast Pre-Score - Test 02 (11 Nov 2020)}, paper=1, qno=9, type=multi, subtopic={Area}, answer={A,B,C,D}, answer_src=key, difficulty=8, diff_parts={K2 S2 I2 C2 F1}, diff_note={concepts: triangle from two medians, area | steps: centroid relations; side lengths; area; chord of circumcircle}]
In a triangle ABC length of medians AD and CE are 18 and 27 respectively. CE is extended to intersect its circumcircle of ABC at F. If AB = 24 and G is centroid of the triangle, then-
\begin{qoptions}
\opt{A}{b = $3\sqrt{70}$}
\opt{B}{a = $6\sqrt{31}$}
\opt{C}{$\Delta$AGE = $\dfrac{27}{4}\sqrt{55}$}
\opt{D}{EF = $\dfrac{16}{3}$}
\end{qoptions}
\end{bankq}
\begin{bankq}[id=EPS02-P1-MATH-S2-Q10, source={Enthusiast Pre-Score - Test 02 (11 Nov 2020)}, paper=1, qno=10, type=multi, subtopic={Circles of a triangle}, answer={A,C}, answer_src=key, difficulty=8, diff_parts={K2 S2 I2 C2 F1}, diff_note={concepts: isosceles triangle with given cosine, circles of a triangle | steps: find the sides; area; circumradius}]
In an isosceles $\Delta$ABC with integral sides, AB = 7 cm and $\cos A=\dfrac27$, then which of the following is/are true
\begin{qoptions}
\opt{A}{Area of $\Delta$ABC is $6\sqrt5$ cm$^2$}
\opt{B}{Area of $\Delta$ABC is $\dfrac{1029}{32}\sqrt5$ cm$^2$}
\opt{C}{Circumradius of $\Delta$ABC is $\dfrac{49\sqrt5}{30}$ cm}
\opt{D}{Circumradius of $\Delta$ABC is $\dfrac{49\sqrt5}{12}$ cm}
\end{qoptions}
\end{bankq}
\begin{bankq}[id=EPS02-P1-MATH-S2-Q11, source={Enthusiast Pre-Score - Test 02 (11 Nov 2020)}, paper=1, qno=11, type=multi, subtopic={Sine and cosine rules}, answer={A,B,C}, answer_src=key, difficulty=6, diff_parts={K2 S1 I2 C1 F1}, diff_note={concepts: progression of sides and trigonometric quantities, sine and cosine rules | steps: relations for a b c in AP; check each quantity}]
In $\Delta$ABC, if a,b,c are in A.P., then which of the following are always in A.P. (symbols have their usual meaning)
\begin{qoptions}
\opt{A}{sinA, sinB, sinC}
\opt{B}{$\dfrac1{r_1},\dfrac1{r_2},\dfrac1{r_3}$}
\opt{C}{$\cot\dfrac A2,\cot\dfrac B2,\cot\dfrac C2$}
\opt{D}{Distance of circumcentre from sides BC, CA and AB respectively.}
\end{qoptions}
\end{bankq}
\begin{bankq}[id=EPS02-P1-MATH-S3-Q04, source={Enthusiast Pre-Score - Test 02 (11 Nov 2020)}, paper=1, qno=4, type=numerical, subtopic={Sine and cosine rules}, answer={1.50}, answer_src=key, difficulty=8, diff_parts={K2 S2 I2 C2 F1}, diff_note={concepts: perpendicular median and angle bisector, sine and cosine rules | steps: geometry of the figure; relations; ratio}]
In a $\Delta$ABC, if the internal angle bisector of angle A meeting BC at L and median BE are perpendicular, then $\dfrac{BE\cot\frac A2}{AL}$ is equal to
\end{bankq}
\begin{bankq}[id=EPS02-P1-MATH-S3-Q05, source={Enthusiast Pre-Score - Test 02 (11 Nov 2020)}, paper=1, qno=5, type=numerical, subtopic={Circles of a triangle}, answer={5.00}, answer_src=key, difficulty=8, diff_parts={K2 S2 I2 C2 F1}, diff_note={concepts: cyclic quadrilateral with given area, circles of a triangle | steps: chord relations; angles; area; perimeter}]
A cyclic quadrilateral ABCD of area $\dfrac{3\sqrt3}{4}$ is inscribed in a unit circle. If AB = 1, BD = $\sqrt3$, then perimeter of cyclic quadrilateral is
\end{bankq}
\begin{bankq}[id=EPS02-P2-MATH-S1-Q05, source={Enthusiast Pre-Score - Test 02 (11 Nov 2020)}, paper=2, qno=5, type=multi, subtopic={Circles of a triangle}, answer={A,C}, answer_src=key, difficulty=8, diff_parts={K2 S2 I2 C2 F1}, diff_note={concepts: isosceles triangle with equal circumradius, circles of a triangle | steps: angles; sine sum; inradius ratio}]
Let the radius of circumcircle of an isosceles triangle ABC is equal to AB = AC. If $P=\sin^2A+\sin^2B$ and $Q=\dfrac rR$, then (where r and R are inradius and circumradius of $\Delta$ABC respectively)-
\begin{qoptions}
\opt{A}{P = 1}
\opt{B}{$P=\dfrac12$}
\opt{C}{$Q=\dfrac{2\sqrt3-3}{2}$}
\opt{D}{$Q=\dfrac{3\sqrt3-2}{2}$}
\end{qoptions}
\end{bankq}
\begin{bankq}[id=EPS02-P2-MATH-S1-Q06, source={Enthusiast Pre-Score - Test 02 (11 Nov 2020)}, paper=2, qno=6, type=multi, subtopic={Sine and cosine rules}, answer={B,C,D}, answer_src=key, difficulty=8, diff_parts={K2 S2 I2 C2 F1}, diff_note={concepts: triangle with sides in AP, sine and cosine rules | steps: half angle relations; tangent rule; products}]
If in $\Delta$ABC, sides a,b,c are in AP and $\cos C=\dfrac13$, then (notations have their usual meanings)-
\begin{qoptions}
\opt{A}{$\Delta$ABC is obtuse angle triangle}
\opt{B}{$\tan\left(\dfrac{B-A}{2}\right)=\dfrac{\sqrt2}{19}$}
\opt{C}{2Rr = 33 for a = 9}
\opt{D}{$\sin\dfrac A2\sin\dfrac B2\sin\dfrac C2=\dfrac{4}{33}$}
\end{qoptions}
\end{bankq}
\begin{bankq}[id=EPS02-P2-MATH-S1-Q07, source={Enthusiast Pre-Score - Test 02 (11 Nov 2020)}, paper=2, qno=7, type=multi, subtopic={Sine and cosine rules}, answer={A,B,D}, answer_src=key, difficulty=9, diff_parts={K3 S2 I2 C2 F1}, diff_note={concepts: triangle from a quadratic in sines, sine and cosine rules, complete squares | steps: complete squares; find angles; evaluate sums}]
In $\Delta$ABC, $3\sin^2B+4\sin B\sin A+16\sin^2A-8\sin B-20\sin A+9=0$, then -
\begin{qoptions}
\opt{A}{$\sin^2A+\sin^2B+\sin^2C=2$}
\opt{B}{$\cos2A+\cos2B+\cos2C=-1$}
\opt{C}{$\cot A+\cot B+\cot C=\sqrt3$}
\opt{D}{$\sin2A+\sin2B+\sin2C=\sqrt3$}
\end{qoptions}
\end{bankq}
\begin{bankq}[id=EPS02-P2-MATH-S3-Q02, source={Enthusiast Pre-Score - Test 02 (11 Nov 2020)}, paper=2, qno=2, type=numerical, subtopic={Circles of a triangle}, answer={9.00}, answer_src=key, difficulty=8, diff_parts={K2 S2 I2 C2 F1}, diff_note={concepts: triangle condition and maximum area, circles of a triangle | steps: simplify b cot B + c cot C; relate to r and R; maximise area}]
With usual notations, if in a $\Delta$ABC, a = 6 \& bcotB + ccotC is equal to twice of sum of its inradius \& its circumradius, then its maximum area is
\end{bankq}
\begin{bankq}[id=EPS02-P2-MATH-S3-Q03, source={Enthusiast Pre-Score - Test 02 (11 Nov 2020)}, paper=2, qno=3, type=numerical, subtopic={Circles of a triangle}, answer={4.00}, answer_src=key, difficulty=7, diff_parts={K2 S2 I2 C1 F1}, diff_note={concepts: sum identity in a triangle, circles of a triangle | steps: simplify the sum; relate to the product of half angles; solve R}]
In $\Delta$ABC, if $\dfrac{\sum(b+c)\cos A}{4\cos\frac A2\cos\frac B2\cos\frac C2}=8$, then the value of circumradius of $\Delta$ABC, is
\end{bankq}
\begin{bankq}[id=EPS02-P2-MATH-S3-Q05, source={Enthusiast Pre-Score - Test 02 (11 Nov 2020)}, paper=2, qno=5, type=numerical, subtopic={Circles of a triangle}, answer={0.60}, answer_src=key, difficulty=8, diff_parts={K2 S2 I2 C2 F1}, diff_note={concepts: perpendicular from circumcentre, circles of a triangle | steps: angle of AD; relation to B-C; solve R}]
In acute angled triangle ABC with usual notations, perpendicular from circumcentre upon side BC intersects circumcircle at D, AD = $\sqrt{10}$, $\sin(B-C)=\dfrac35$, then $\dfrac1R$ is (where R is circumradius of triangle ABC)
\end{bankq}
\begin{bankq}[id=EPS03-P2-MATH-S3-Q07, source={Enthusiast Pre-Score - Test 03 (23 Nov 2020)}, paper=2, qno=7, type=numerical, subtopic={Circles of a triangle}, answer={5.00}, answer_src=key, difficulty=8, diff_parts={K2 S2 I2 C2 F1}, diff_note={concepts: maximum of sum of distances from a point on the arc, circles of a triangle | steps: Ptolemy; parameter; maximise}]
Let M be any point M on the arc AB of circumcircle of right angled triangle ABC and the maximum value of AM+BM+CM is $\sqrt{29N}$, where $a=3$, $b=4$, $c=5$. Then N is $\ldots$
\end{bankq}

\section*{TOAS papers}
\begin{bankq}[id=TOAS01-P2-MATH-S1-Q01, source={TOAS - Test 01 (30 Jan 2022)}, paper=2, qno=1, type=single, subtopic={Area}, answer={C}, answer_src=key, difficulty=2, diff_parts={K1 S1 I1 C0 F0}, diff_note={concepts: areas in a triangle divided by cevians | steps: ratio of areas of triangles sharing a height; solve for x | I: area ratio technique}]
A triangle is divided into 4 pieces and the area of the 4 regions are given inside the corresponding regions, then value of $x$ is -
\bfig{79.6}{TOAS01_p2_math_s1_q01_1.png}
\begin{qoptions}
\opt{A}{18}
\opt{B}{15}
\opt{C}{22}
\opt{D}{12.5}
\end{qoptions}
\end{bankq}
\begin{bankq}[id=TOAS05-P1-MATH-S1-Q02, source={TOAS - Test 05 (03 Apr 2022)}, paper=1, qno=2, type=multi, subtopic={Sine and cosine rules}, answer={A,B,C}, answer_src=key, difficulty=6, diff_parts={K2 S3 I1 C0 F1}, diff_note={concepts: sine rule, angle bisector | steps: C = 2B; AD bisects A; AC = BD; sine rule relation; solve the cosines | I: relate AC = BD to the angle bisector ratio}, also_in={TAAS - Test 01 (18 Jun 2023) P1 Q5; TNAS - Test 04 (18 Jun 2023) P1 Q5}]
$ABC$ is a triangle in which $\angle C=2\angle B$. $D$ is a point on $BC$ such that $AD$ bisects $\angle BAC$ and $AC=BD$. Then which of the following is/are true
\begin{qoptions}
\opt{A}{$\cos B+\cos C=\dfrac{\sqrt5}{2}$}
\opt{B}{$\cos B-\cos C=\dfrac12$}
\opt{C}{$\dfrac{\sin A}{\sin C}=1$}
\opt{D}{$\dfrac{\sin A}{\sin C}=\dfrac{\sqrt5+1}{4}$}
\end{qoptions}
\end{bankq}
\begin{bankq}[id=TOAS05-P1-MATH-S1-Q03, source={TOAS - Test 05 (03 Apr 2022)}, paper=1, qno=3, type=multi, subtopic={Sine and cosine rules}, answer={A,C}, answer_src=key, difficulty=5, diff_parts={K2 S2 I1 C0 F1}, diff_note={concepts: trigonometry, coordinate geometry from a figure | steps: set up lengths from the figure; apply tangent and sine relations; solve for tan alpha | I: choose the right auxiliary angle}, also_in={TNAS - Test 04 (18 Jun 2023) P2 Q8}]
For the figure given below if $\tan\alpha=\dfrac mn$, where $m$ and $n$ are relatively prime positive integers then
\bfig{84.0}{TOAS05_p1_math_s1_q03_1.png}
\begin{qoptions}
\opt{A}{$m+n=463$}
\opt{B}{$n-m=122$}
\opt{C}{$2m-n=41$}
\opt{D}{$m+n=465$}
\end{qoptions}
\end{bankq}
\begin{bankq}[id=TOAS05-P1-MATH-S1-Q07, source={TOAS - Test 05 (03 Apr 2022)}, paper=1, qno=7, type=multi, subtopic={Sine and cosine rules}, answer={A}, answer_src=key, difficulty=3, diff_parts={K1 S1 I1 C0 F1}, diff_note={concepts: angle chasing in a triangle | steps: exterior angle relations at D, E, F; combine with the equilateral triangle | I: exterior angle sums}]
In the diagram shown, $AB$ and $AC$ are equal sides of an isosceles triangle $ABC$. The equilateral triangle $DEF$ is inscribed. If $\angle ADF$ be $\alpha$, $\angle BED$ be $\beta$ and $\angle CFE$ be $\gamma$, then \textbf{CORRECT} relation is/are
\bfig{105.3}{TOAS05_p1_math_s1_q07_1.png}
\begin{qoptions}
\opt{A}{$\beta=\dfrac{\alpha+\gamma}2$}
\opt{B}{$\alpha=\dfrac{\beta+\gamma}2$}
\opt{C}{$\beta=\dfrac{\alpha-\gamma}2$}
\opt{D}{$\alpha=\dfrac{\beta-\gamma}2$}
\end{qoptions}
\end{bankq}
\begin{bankq}[id=TOAS05-P1-MATH-S3-Q02, source={TOAS - Test 05 (03 Apr 2022)}, paper=1, qno=2, type=numerical, subtopic={Sine and cosine rules}, answer={8.33}, answer_src=key, difficulty=6, diff_parts={K2 S3 I1 C0 F1}, diff_note={concepts: medians and centroid, sine rule | steps: AG = 10/3; angle at G in triangle AGC; CG by the sine rule; area of AGC; multiply by 3 | I: work in triangle AGC}, also_in={TAAS - Test 01 (18 Jun 2023) P1 Q4; TNAS - Test 04 (18 Jun 2023) P1 Q4}]
In a triangle $ABC$, medians $AD$ and $CE$ are drawn. If $AD=5$, $\angle DAC=\dfrac\pi8$ and $\angle ACE=\dfrac\pi4$ then the area of the triangle $ABC$ is equal to
\end{bankq}
\begin{bankq}[id=TOAS05-P2-MATH-S1-Q20, source={TOAS - Test 05 (03 Apr 2022)}, paper=2, qno=20, type=single, subtopic={Area}, answer={B}, answer_src=key, difficulty=2, diff_parts={K1 S2 I0 C0 F0}, diff_note={concepts: area of a trapezium | steps: area formula; relation a = 2b; solve for c}, also_in={TMAS - Test 01 (09 Jul 2023) P1 Q10; TNAS - Test 03 (28 May 2023) P1 Q10}]
As shown in the figure, the area of the trapezium $ABCD$ is $\dfrac{81}2$, where $a=2b$, then $c$ is equal to -
\bfig{74.9}{TOAS05_p2_math_s1_q20_1.png}
\begin{qoptions}
\opt{A}{$\dfrac92$}
\opt{B}{$3\sqrt5$}
\opt{C}{$\sqrt{30}$}
\opt{D}{none of these}
\end{qoptions}
\end{bankq}
\begin{bankq}[id=TOAS05-P2-MATH-S2-Q08, source={TOAS - Test 05 (03 Apr 2022)}, paper=2, qno=8, type=numerical, subtopic={Sine and cosine rules}, answer={786}, answer_src=solved, working={$a^2+c^2-b^2=ac$ so $\cos B=\frac12$, $\sin^2(\pi/4)=\frac12$, $1572\cdot\frac12=786$}, difficulty=3, diff_parts={K1 S2 I0 C0 F1}, diff_note={concepts: cosine rule | steps: rearrange to a\^{}2 + c\^{}2 - b\^{}2 = ac; cos B = 1/2; evaluate}]
If in a triangle $ABC$ with usual notation, $a(a-c)=b^2-c^2$, then $1572(\sin^2(\pi\cos^2B))$ is equal to
\end{bankq}
\begin{bankq}[id=TOAS05-P2-MATH-S2-Q09, source={TOAS - Test 05 (03 Apr 2022)}, paper=2, qno=9, type=numerical, subtopic={Circles of a triangle}, answer={2}, answer_src=solved, working={$\frac1r=\frac14+\frac16+\frac18=\frac{13}{24}$, $r=\frac{24}{13}$; $2q-p=26-24=2$}, difficulty=2, diff_parts={K1 S1 I0 C0 F1}, diff_note={concepts: exradii and inradius | steps: 1/r = 1/r1 + 1/r2 + 1/r3}]
If $r_1=4$, $r_2=6$, $r_3=8$ for $\Delta ABC$ and $r=\dfrac pq$ (where $p$ \& $q$ are relatively prime numbers), then value of $(2q-p)$ is (usual notation is followed)
\end{bankq}
\begin{bankq}[id=TOAS05-P2-MATH-S2-Q10, source={TOAS - Test 05 (03 Apr 2022)}, paper=2, qno=10, type=numerical, subtopic={Circles of a triangle}, answer={7}, answer_src=key, difficulty=3, diff_parts={K1 S1 I1 C0 F1}, diff_note={concepts: half angle formulas | steps: cot(A/2) cot(B/2) = s/(s - c); use a + b = 4c/3 | I: product of half-angle cotangents}]
If in a $\Delta ABC$ with usual notations, $3a+3b=4c$, then $\cot\dfrac A2\cdot\cot\dfrac B2$ is
\end{bankq}

\section*{TAAS papers}
\begin{bankq}[id=TAAS01-P1-MATH-S2-Q12, source={TAAS - Test 01 (18 Jun 2023)}, paper=1, qno=12, type=multi, subtopic={Sine and cosine rules}, answer={A,B,D}, answer_src=key, difficulty=4, diff_parts={K1 S2 I1 C0 F1}, diff_note={concepts: solving a triangle with the cosine rule | steps: largest angle; tangent of half the angle difference; median length; angle bisector length}, also_in={TNAS - Test 04 (18 Jun 2023) P1 Q12}]
If the sides $a,b,c$ of a triangle are 3, 5 and 7 respectively then which of the following is/are true ?
\begin{qoptions}
\opt{A}{$\tan\left(\dfrac{A-B}2\right)=\dfrac{-1}{4\sqrt3}$}
\opt{B}{largest angle is $120^\circ$}
\opt{C}{length of median through vertex $A$ is $\dfrac{\sqrt{139}}4$}
\opt{D}{length of angle bisector through $C$ is $\dfrac{15}8$}
\end{qoptions}
\end{bankq}

\section*{TOPS papers}
\begin{bankq}[id=TOPS03-P2-MATH-S1-Q02, source={TOPS - Test 03 (20 Nov 2022)}, paper=2, qno=2, type=multi, subtopic={Circles of a triangle}, answer={A,B,D}, answer_src=key, difficulty=7, diff_parts={K2 S2 I2 C1 F1}, diff_note={concepts: arcs and inscribed angles, area of a triangle | steps: angles from the arcs; check acute and A.P.; radius of the circle; area | triangle on a circle}]
Let $A,B$ and $C$ are three points on a circle such that arc $AB=6$, arc $BC=8$, arc $CA=10$, then
\begin{qoptions}
\opt{A}{$\Delta ABC$ is acute angled}
\opt{B}{$\angle B,\angle A,\angle C$ are in A.P}
\opt{C}{$\angle A,\angle B,\angle C$ are in A.P}
\opt{D}{Area of $\Delta ABC$ is $\dfrac{36\left(3+\sqrt3\right)}{\pi^2}$}
\end{qoptions}
\end{bankq}
\begin{bankq}[id=TOPS03-P2-MATH-S3-Q05, source={TOPS - Test 03 (20 Nov 2022)}, paper=2, qno=5, type=numerical, subtopic={Sine and cosine rules}, answer={1}, answer_src=key, difficulty=10, diff_parts={K3 S3 I2 C2 F1}, diff_note={concepts: chord length, circumradius, area of a triangle | steps: chord relations; area formula; angle at the centre; sum of squares; radius | circle and chords}]
Let a circle with centre $O$ and radius $k$. Three points $P$, $Q$, $R$ are chosen on a circle such that if length of chord $PQ$, $QR$, $RP$ are $x$, $y$ \& $z$ respectively, $x^2+y^2+z^2=5$, $\angle QOR=60$ and Area of $\Delta PQR$ is $\dfrac{\sqrt3}4$ unit, then the value of radius $k$ of the circle is equal to
\end{bankq}
\begin{bankpassage}[id=TOPS04-P1-MATH-S2-G9-10, source={TOPS - Test 04 (18 Dec 2022)}, questions={9-10}]
In a triangle $ABC$ if $r_1=2r_2=3r_3$ then (where $r_1,r_2,r_3$ are ex-radius of triangle $ABC$)
\end{bankpassage}
\begin{bankq}[id=TOPS04-P1-MATH-S2-Q09, source={TOPS - Test 04 (18 Dec 2022)}, paper=1, qno=9, type=single, subtopic={Circles of a triangle}, answer={B}, answer_src=key, difficulty=5, diff_parts={K2 S2 I1 C1 F0}, diff_note={concepts: exradii, half angle tangent | steps: ratio of the exradii; relation of s minus a; tangent ratio}, group=TOPS04-P1-MATH-S2-G9-10]
The ratio of $\tan\dfrac A2:\tan\dfrac B2:\tan\dfrac C2$ is
\begin{qoptions}
\opt{A}{$6:5:4$}
\opt{B}{$6:3:2$}
\opt{C}{$4:5:6$}
\opt{D}{$3:2:6$}
\end{qoptions}
\end{bankq}
\begin{bankq}[id=TOPS04-P1-MATH-S2-Q10, source={TOPS - Test 04 (18 Dec 2022)}, paper=1, qno=10, type=single, subtopic={Circles of a triangle}, answer={C}, answer_src=key, difficulty=5, diff_parts={K2 S2 I1 C1 F0}, diff_note={concepts: exradii, sides | steps: s minus each side ratio; sides from s; ratio}, group=TOPS04-P1-MATH-S2-G9-10]
The ratio of $a:b:c$ is (where $a,b,c$ are sides opposite to angles $A,B,C$ of triangle $ABC$) -
\begin{qoptions}
\opt{A}{$3:4:5$}
\opt{B}{$4:3:5$}
\opt{C}{$5:4:3$}
\opt{D}{$5:3:4$}
\end{qoptions}
\end{bankq}
\begin{bankq}[id=TOPS04-P2-MATH-S2-Q01, source={TOPS - Test 04 (18 Dec 2022)}, paper=2, qno=1, type=numerical, subtopic={Circles of a triangle}, answer={1.00}, answer_src=key, difficulty=6, diff_parts={K1 S2 I2 C1 F1}, diff_note={concepts: half angle tangents | steps: relation of s and a; tangent product; value}]
In a $\Delta ABC$, with usual notations, $4b+4c=5a$, then $9\left(\tan\dfrac C2\tan\dfrac B2\right)$ is
\end{bankq}
\begin{bankq}[id=TOPS04-P2-MATH-S2-Q05, source={TOPS - Test 04 (18 Dec 2022)}, paper=2, qno=5, type=numerical, subtopic={Circles of a triangle}, answer={25}, answer_src=solved, working={Sum of coefficients $=\left(\tfrac1{r_1}+\tfrac1{r_2}+\tfrac1{r_3}\right)^{20}=\left(\tfrac12\right)^{20}$. Then $50\times2^{-20}\times2^{19}=25$}, difficulty=6, diff_parts={K2 S2 I1 C1 F1}, diff_note={concepts: inradius and exradii, coefficients | steps: relation of reciprocals; sum of coefficients; evaluate}]
In $\Delta ABC$ with usual notations, let the length of inradius is 2 and sum of coefficient in the expansion of $\left(\left(\dfrac1{r_1}+\dfrac1{r_2}\right)x+\dfrac y{r_3}\right)^{20}$ is $\lambda$ then $50(\lambda\times2^{19})$ is (if expansion is written in powers of $x$ and $y$)
\end{bankq}
\begin{bankq}[id=TOPS05-P2-MATH-S2-Q08, source={TOPS - Test 05 (14 Jan 2023)}, paper=2, qno=8, type=multi, subtopic={Sine and cosine rules}, answer={C,D}, answer_src=key, difficulty=9, diff_parts={K2 S3 I2 C2 F1}, diff_note={concepts: triangle, angle bisector | steps: angles; bisector length; sides by the sine rule; circumradius; area}]
The triangle $ABC$ has $\angle A=60^\circ$, $\angle B=45^\circ$. The bisector of $\angle A$ meets the side $BC$ at $T$ where $AT=24$. Then which of the following is correct?
\begin{qoptions}
\opt{A}{length of $BC=18\sqrt2$}
\opt{B}{length of $AC<AT$}
\opt{C}{radius of the circumcircle is $12\sqrt2$}
\opt{D}{area of the triangle is $72\left(3+\sqrt3\right)$ sq. units.}
\end{qoptions}
\end{bankq}
\begin{bankq}[id=TOPS05-P2-MATH-S3-Q07, source={TOPS - Test 05 (14 Jan 2023)}, paper=2, qno=7, type=numerical, subtopic={Sine and cosine rules}, answer={5}, answer_src=key, difficulty=7, diff_parts={K2 S2 I2 C1 F1}, diff_note={concepts: triangle, sine rule | steps: tan C from A and B; sin C; circumradius; evaluate}]
In a $\Delta ABC$, $\tan A=2$, $\tan B=\dfrac32$ and $c=2\sqrt{65}$. If circumradius of the triangle is $k$, then value of $(7k/13)$ is
\end{bankq}

\section*{TNAS papers}
\begin{bankq}[id=TNAS01-P1-MATH-S1-Q15, source={TNAS - Test 01 (05 Feb 2023)}, paper=1, qno=15, type=single, subtopic={Circles of a triangle}, answer={C}, answer_src=key, difficulty=5, diff_parts={K1 S2 I2 C1 F0}, diff_note={concepts: triangle relations | steps: use A minus B; R and r relation; cosine of C}]
In triangle ABC, if $A-B=90^\circ$ and $R=4r$, where $R$ and $r$ have their usual meanings, then $\cos C$ is equal to
\begin{qoptions}
\opt{A}{$\dfrac23$}
\opt{B}{$\dfrac78$}
\opt{C}{$\dfrac34$}
\opt{D}{$\dfrac{\sqrt3}{2}$}
\end{qoptions}
\end{bankq}
\begin{bankq}[id=TNAS01-P1-MATH-S2-Q05, source={TNAS - Test 01 (05 Feb 2023)}, paper=1, qno=5, type=numerical, subtopic={Sine and cosine rules}, answer={7}, answer_src=key, difficulty=6, diff_parts={K1 S2 I2 C1 F1}, diff_note={concepts: triangle bisector lengths | steps: sides from sine rule; angles; bisector formula; ratio}]
The ratio of the length of the sides BC and AC of $\triangle$ABC to the radius of circumscribed circle are equal to $2$ and $\dfrac32$ respectively. If the ratio of the length of the bisectors of the interior angles B and C is $\dfrac{\alpha(\sqrt\alpha-1)}{\beta\sqrt\gamma}$ (where $\alpha,\beta,\gamma\in N$), then $\alpha$ is
\end{bankq}

\section*{Star Review (Other) papers}
\begin{bankq}[id=STAR01-P1-MATH-S1-Q04, source={Star Review (Other) - Test 01 (12 Sep 2022)}, paper=1, qno=4, type=multi, subtopic={Sine and cosine rules}, answer={B,C,D}, answer_src=key, difficulty=8, diff_parts={K1 S3 I2 C2 F1}, diff_note={concepts: triangle condition and angle bound | steps: simplify the given relation; find the angle C; series evaluation; maximise the ratio; check}]
In $\Delta ABC$, if $\dfrac1{a+c}+\dfrac1{b+c}=\dfrac3{a+b+c}$ then which of the following hold good ?
\begin{qoptions}
\opt{A}{$\displaystyle\sum_{r=1}^{\infty}\cos^{2r}(A+B)=1$}
\opt{B}{$\displaystyle\sum_{r=1}^{\infty}\cos^{2r}(A+B)=\dfrac13$}
\opt{C}{$\dfrac{\cot C}{\cot A+\cot B}$ is maximum if $\Delta ABC$ is equilateral}
\opt{D}{$\dfrac{\cot C}{\cot A+\cot B}$ is maximum if $\Delta ABC$ is isosceles}
\end{qoptions}
\end{bankq}
\begin{bankpassage}[id=STAR01-P1-MATH-S3-G1-2, source={Star Review (Other) - Test 01 (12 Sep 2022)}, questions={1-2}]
In $\Delta ABC$, if $a=4$, $b=3$, $c=2$ then
\end{bankpassage}
\begin{bankq}[id=STAR01-P1-MATH-S3-Q01, source={Star Review (Other) - Test 01 (12 Sep 2022)}, paper=1, qno=1, type=numerical, subtopic={Sine and cosine rules}, answer={2.25}, answer_src=key, difficulty=5, diff_parts={K1 S2 I1 C1 F1}, diff_note={concepts: triangle with given sides | steps: cosine rule for B; cosine rule for C; evaluate}, group=STAR01-P1-MATH-S3-G1-2]
The value of $2\cos B+\cos C$ is
\end{bankq}
\begin{bankq}[id=STAR01-P1-MATH-S3-Q02, source={Star Review (Other) - Test 01 (12 Sep 2022)}, paper=1, qno=2, type=numerical, subtopic={Sine and cosine rules}, answer={3.50}, answer_src=key, difficulty=5, diff_parts={K1 S2 I1 C1 F1}, diff_note={concepts: ratio of sines in a triangle | steps: sine rule; angle relation; ratio}, group=STAR01-P1-MATH-S3-G1-2]
The value of $\dfrac{\sin A+\sin B}{\sin(A+B)}$ is
\end{bankq}
\begin{bankq}[id=STAR01-P2-MATH-S1-Q10, source={Star Review (Other) - Test 01 (12 Sep 2022)}, paper=2, qno=10, type=single, subtopic={Sine and cosine rules}, answer={D}, answer_src=key, difficulty=4, diff_parts={K1 S2 I1 C1 F0}, diff_note={concepts: third side from the roots of a quadratic | steps: sum and product of roots; cosine rule; evaluate}]
Two sides `a' and `b' of triangle ABC are given by the roots of the equation $x^2-4x+1=0$ and the included angle between them is $\dfrac\pi3$, then the square of the third side is -
\begin{qoptions}
\opt{A}{$\dfrac14$}
\opt{B}{4}
\opt{C}{11}
\opt{D}{13}
\end{qoptions}
\end{bankq}
\begin{bankq}[id=STAR01-P2-MATH-S2-Q01, source={Star Review (Other) - Test 01 (12 Sep 2022)}, paper=2, qno=1, type=numerical, subtopic={Sine and cosine rules}, answer={2}, answer_src=key, difficulty=6, diff_parts={K1 S2 I2 C1 F1}, diff_note={concepts: triangle with a relation between half angles | steps: use the sum formula; compare the terms; find the angle; scale}]
In $\Delta ABC$ if $2\cos\dfrac B2\cos\dfrac C2=\dfrac12+\dfrac{(b+c)}a\sin\dfrac A2$ and $A$ is measured in radian then $\dfrac{6A}\pi$ is
\end{bankq}
\begin{bankq}[id=STAR01-P2-MATH-S2-Q06, source={Star Review (Other) - Test 01 (12 Sep 2022)}, paper=2, qno=6, type=numerical, subtopic={Sine and cosine rules}, answer={5}, answer_src=solved, working={$\cos\theta=\dfrac{(2x+1)^2+(x^2-1)^2-(x^2+x+1)^2}{2(2x+1)(x^2-1)}=-\tfrac12$, so $\theta=\tfrac{2\pi}3$ and $a+b=5$}, difficulty=8, diff_parts={K1 S3 I2 C2 F1}, diff_note={concepts: greatest angle of a triangle with polynomial sides | steps: largest side; cosine rule; simplify; angle; sum}]
The sides of a triangle are $x^2+x+1$, $2x+1$ and $x^2-1$. Let $\theta$ be the greatest angle of triangle. If $\theta=\dfrac{a\pi}b$, where $a$, $b$ are coprime numbers, then value of $(a+b)$ is
\end{bankq}
\begin{bankq}[id=STAR02-P1-MATH-S2-Q10, source={Star Review (Other) - Test 02 (10 Oct 2022)}, paper=1, qno=10, type=single, subtopic={Sine and cosine rules}, answer={A}, answer_src=key, difficulty=6, diff_parts={K1 S3 I2 C1 F0}, diff_note={concepts: triangle with two sides and the included angle | steps: cosine rule; circumradius; angle bisector length; area; match}]
A triangle $ABC$ is shown in the figure with $\angle C=60^\circ$, $a=5$ \& $b=4$.\par
\bfig{75.1}{STAR02_p1_math_s2_q10_1.png}
Match List-I with List-II and select the correct answer using the code given below the list.\par
{\small\setlength{\tabcolsep}{2pt}\begin{tabular}{|c|p{0.4\linewidth}|c|c|}
\hline
\multicolumn{2}{|c|}{\textbf{List-I}} & \multicolumn{2}{c|}{\textbf{List-II}}\\ \hline
(I) & The value of the side opposite to angle $\angle C$ is & (P) & $5\sqrt3$\\ \hline
(II) & If $R$ is the circumradius of the triangle $ABC$, the value of $R^2$ is & (Q) & $7$\\ \hline
(III) & Length of internal angle bisector of $C$ is & (R) & $\sqrt{21}$\\ \hline
(IV) & Area of triangle $ABC$ is & (S) & $\dfrac{20}{3\sqrt3}$\\ \hline
 & & (T) & $3\sqrt5$\\ \hline
\end{tabular}}\par
\begin{qoptions}
\opt{A}{$I\to R;\allowbreak II\to Q;\allowbreak III\to S;\allowbreak IV\to P$}
\opt{B}{$I\to Q;\allowbreak II\to R;\allowbreak III\to S;\allowbreak IV\to P$}
\opt{C}{$I\to P;\allowbreak II\to Q;\allowbreak III\to R;\allowbreak IV\to S$}
\opt{D}{$I\to R;\allowbreak II\to S;\allowbreak III\to P;\allowbreak IV\to Q$}
\end{qoptions}
\end{bankq}

\section*{Nurture (General) papers}
\begin{bankq}[id=NUR01-P2-MATH-S1-Q05, source={Nurture (General) - Test 01 (12 Jun 2021)}, paper=2, qno=5, type=single, subtopic={Sine and cosine rules}, answer={C}, answer_src=solved, working={Sum of squares of the two cevians is 5h\^{}2/9; h=3/sqrt5}, difficulty=6, diff_parts={K2 S2 I2 C1 F0}, diff_note={concepts: cevians to the trisection points of a hypotenuse, sine and cosine rules | steps: set coordinates; sum of squares; solve for the hypotenuse}]
Line segments drawn from vertex opposite to the hypotenuse of a right angle triangle to the points trisecting the hypotenuse, have lengths $\sin x$ and $\cos x$ $\left(0<x<\dfrac\pi2\right)$. Then length of hypotenuse is
\begin{qoptions}
\opt{A}{$\dfrac13$}
\opt{B}{$\dfrac32$}
\opt{C}{$\dfrac{3\sqrt5}5$}
\opt{D}{can't be determined}
\end{qoptions}
\end{bankq}
\begin{bankq}[id=NUR04-P1-MATH-S3-Q18, source={Nurture (General) - Test 04 (11 Feb 2024)}, paper=1, qno=18, type=numerical, subtopic={Area}, answer={24}, answer_src=solved, working={Triangle 3-4-5 with a point at angles 90, 120, 150: 2xy+sqrt3 yz+xz=4x area=24}, difficulty=8, diff_parts={K2 S2 I2 C2 F1}, diff_note={concepts: system of equations as a geometric triangle, area | steps: interpret as a triangle with a point; angles; area; relation}]
Three positive reals x, y, z satisfy\par
$x^2+y^2=3^2$\par
$y^2+yz+z^2=4^2$\par
$x^2+\sqrt3xz+z^2=5^2$\par
Find the value of $2xy+xz+\sqrt3yz$
\end{bankq}
\end{document}
